\begin{abstract}
Extracorporeal membrane oxygenation (ECMO) is an essential life-supporting modality for COVID-19 patients who are refractory to conventional therapies. However, the proper treatment decision  has been the subject of significant debate and it remains controversial about who benefits from this scarcely available and technically complex treatment option. To support clinical decisions, it is a critical need to predict the treatment need and the potential treatment and no-treatment responses. Targeting this clinical challenge, we propose Treatment Variational AutoEncoder (TVAE), a novel approach for individualized treatment analysis. TVAE is specifically designed to address the modeling challenges like ECMO with strong treatment selection bias and scarce treatment cases. TVAE conceptualizes the treatment decision as a multi-scale problem. We model a patient's potential treatment assignment and the factual and counterfactual outcomes as part of their intrinsic characteristics that can be represented by a deep latent variable model. The factual and counterfactual prediction errors are alleviated via a reconstruction regularization scheme together with semi-supervision, and the selection bias and the scarcity of treatment cases are mitigated by the disentangled and distribution-matched latent space and the label-balancing generative strategy. We evaluate TVAE on two real-world COVID-19 datasets: an international dataset collected from 1651 hospitals across 63 countries, and a institutional dataset collected from 15 hospitals. The results show that TVAE outperforms state-of-the-art treatment effect models in predicting both the propensity scores and factual outcomes on heterogeneous COVID-19 datasets. Additional experiments also show TVAE outperforms the best existing models in individual treatment effect estimation on the synthesized IHDP benchmark dataset.
\end{abstract}

\begin{CCSXML}
<ccs2012>
   <concept>
       <concept_id>10010147.10010257.10010258.10010260.10010271</concept_id>
       <concept_desc>Computing methodologies~Dimensionality reduction and manifold learning</concept_desc>
       <concept_significance>500</concept_significance>
       </concept>
   <concept>
       <concept_id>10010147.10010257.10010282.10011305</concept_id>
       <concept_desc>Computing methodologies~Semi-supervised learning settings</concept_desc>
       <concept_significance>500</concept_significance>
       </concept>
   <concept>
       <concept_id>10010147.10010257.10010293.10010294</concept_id>
       <concept_desc>Computing methodologies~Neural networks</concept_desc>
       <concept_significance>500</concept_significance>
       </concept>
   <concept>
       <concept_id>10010147.10010257.10010293.10010319</concept_id>
       <concept_desc>Computing methodologies~Learning latent representations</concept_desc>
       <concept_significance>500</concept_significance>
       </concept>
   <concept>
       <concept_id>10010405.10010444.10010449</concept_id>
       <concept_desc>Applied computing~Health informatics</concept_desc>
       <concept_significance>500</concept_significance>
       </concept>
   <concept>
       <concept_id>10010405.10010444.10010450</concept_id>
       <concept_desc>Applied computing~Bioinformatics</concept_desc>
       <concept_significance>500</concept_significance>
       </concept>
 </ccs2012>
\end{CCSXML}

\ccsdesc[500]{Computing methodologies~Dimensionality reduction and manifold learning}
\ccsdesc[500]{Computing methodologies~Semi-supervised learning settings}
\ccsdesc[500]{Computing methodologies~Neural networks}
\ccsdesc[500]{Computing methodologies~Learning latent representations}
\ccsdesc[500]{Applied computing~Health informatics}
\ccsdesc[500]{Applied computing~Bioinformatics}


%%
%% Keywords. The author(s) should pick words that accurately describe
%% the work being presented. Separate the keywords with commas.
\keywords{Machine Learning for Healthcare, Causal Inference, Representation Learning, Semi-supervised Learning, Treatment Effect Estimation, COVID Analysis, Generative AI, Deep Latent Variable Models, Variational Autoencoder}
